% GENERATED FILE. Edit canonical lessons, then run npm run generate:books.
% canonical-source-hash: fnv1a64:5856869f4657493e
% canonical-lessons: TA-C504-mutumai, TA-C504-mutiyavar, TA-C504-napakam, TA-C504-manku, TA-C504-curunku

\chapter{Growing Old}
\label{ch:ta-growing-old}
\hlchaptermodality{504}

\begin{chapteropening}
\textbf{By the end of this chapter:} \emph{I can talk about old age, the elderly, memory, fading eyesight and wrinkles.}

Complete the last lesson of chapter 504: I can talk about old age, the elderly, memory, fading eyesight and wrinkles.
\end{chapteropening}

\section[\texorpdfstring{\ta{முதுமை} (mutumai)}{முதுமை (mutumai)}]{\ta{முதுமை} (mutumai) — old age}

\label{lesson:TA-C504-mutumai}



\begin{quote}
Before the new one: say the Tamil for to host, then the Tamil for a farewell.
\end{quote}

\subsection*{You'll want to know: \ta{முதுமை}}
\textbf{\ta{முதுமை}} — \emph{mutumai} — \textquotedblleft{}old age\textquotedblright{}.

\begin{quote}
\ta{முதுமையில்} \ta{தனிமை} \ta{ஒரு} \ta{பெரிய} \ta{பிரச்சினை}.
\end{quote}

\begin{quote}
\emph{mutumaiyil taṉimai oru periya piracciṉai.}
\end{quote}

\begin{quote}
Loneliness is a big problem in old age.
\end{quote}

From \ta{முது} \textquotedblleft{}old\textquotedblright{}; \ta{முதுமையில்}, in old age.

\begin{cousinweb}
\textbf{\ta{முதுமை}} (\emph{mutumai}) is native Dravidian (DEDR 4954), the noun of \ta{முது} \emph{mutu}, \textquotedblleft{}old, ancient\textquotedblright{}.
\end{cousinweb}

\begin{grammarlens}[title={What you've built}]
One more word for describing what you lived and felt.
\end{grammarlens}

\subsection*{Your turn}
\begin{itemize}
\raggedright
  \item \emph{Say it:} \emph{mutumai}
  \item \emph{Say it:} \emph{mutumai}, once more
  \item \emph{Say it:} \emph{mutumaiyil taṉimai oru periya piracciṉai.} — Loneliness is a big problem in old age.
  \item \emph{Recall:} say the Tamil for to host, then the Tamil for a farewell, then say \emph{mutumai} again
\end{itemize}

\begin{tcolorbox}[breakable,colback=teal!4,colframe=teal!35!black,title={Before you move on}]
What does \ta{முதுமை} mean? (\textquotedblleft{}Old age\textquotedblright{}.) Say it once more.
\end{tcolorbox}
\section[\texorpdfstring{\ta{முதியவர்} (mutiyavar)}{முதியவர் (mutiyavar)}]{\ta{முதியவர்} (mutiyavar) — an elderly person}

\label{lesson:TA-C504-mutiyavar}



\begin{quote}
Before the new one: say the Tamil for a farewell, then the Tamil for old age.
\end{quote}

\subsection*{You'll want to know: \ta{முதியவர்}}
\textbf{\ta{முதியவர்}} — \emph{mutiyavar} — \textquotedblleft{}an elderly person\textquotedblright{}.

\begin{quote}
\ta{பேருந்தில்} \ta{நின்ற} \ta{முதியவருக்கு} \ta{நான்} \ta{இடம்} \ta{கொடுத்தேன்}.
\end{quote}

\begin{quote}
\emph{pēruntil niṉṟa mutiyavarukku nāṉ iṭam koṭuttēṉ.}
\end{quote}

\begin{quote}
I gave my place on the bus to an elderly man who was standing.
\end{quote}

Plural: \ta{முதியவர்கள்}, \emph{mutiyavarkaḷ}.

Respectful -\ta{அர்}; plural \ta{முதியவர்கள்}.

\begin{cousinweb}
\textbf{\ta{முதியவர்}} (\emph{mutiyavar}) is native Dravidian, built on \ta{முது} \emph{mutu}, \textquotedblleft{}old\textquotedblright{} (DEDR 4954), with the respectful -\ta{அவர்}.
\end{cousinweb}

\begin{grammarlens}[title={What you've built}]
One more word for describing what you lived and felt.
\end{grammarlens}

\subsection*{Your turn}
\begin{itemize}
\raggedright
  \item \emph{Say it:} \emph{mutiyavar}
  \item \emph{Say it:} \emph{mutiyavar}, once more
  \item \emph{Say it:} \emph{pēruntil niṉṟa mutiyavarukku nāṉ iṭam koṭuttēṉ.} — I gave my place on the bus to an elderly man who was standing.
  \item \emph{Recall:} say the Tamil for a farewell, then the Tamil for old age, then say \emph{mutiyavar} again
\end{itemize}

\begin{tcolorbox}[breakable,colback=teal!4,colframe=teal!35!black,title={Before you move on}]
What does \ta{முதியவர்} mean? (\textquotedblleft{}An elderly person\textquotedblright{}.) Say it once more.
\end{tcolorbox}
\section[\texorpdfstring{\ta{ஞாபகம்} (ñāpakam)}{ஞாபகம் (ñāpakam)}]{\ta{ஞாபகம்} (ñāpakam) — memory; recollection}

\label{lesson:TA-C504-napakam}



\begin{quote}
Before the new one: say the Tamil for old age, then the Tamil for an elderly person.
\end{quote}

\subsection*{You'll want to know: \ta{ஞாபகம்}}
\textbf{\ta{ஞாபகம்}} — \emph{ñāpakam} — \textquotedblleft{}memory; recollection\textquotedblright{}.

\begin{quote}
\ta{தாத்தாவுக்கு} \ta{இப்போது} \ta{ஞாபகம்} \ta{குறைந்துவிட்டது}.
\end{quote}

\begin{quote}
\emph{tāttāvukku ippōtu ñāpakam kuṟaintuviṭṭatu.}
\end{quote}

\begin{quote}
Grandfather’s memory has weakened now.
\end{quote}

\ta{ஞாபகம்} \ta{இருக்கிறது}, I remember.

\begin{cousinweb}
\textbf{\ta{ஞாபகம்}} (\emph{ñāpakam}) is borrowed from Sanskrit \emph{jñāpaka}, \textquotedblleft{}causing to know, informing\textquotedblright{} (Monier-Williams); Tamil writes \emph{jñ} as \ta{ஞ}.
\end{cousinweb}

\begin{grammarlens}[title={What you've built}]
One more word for describing what you lived and felt.
\end{grammarlens}

\subsection*{Your turn}
\begin{itemize}
\raggedright
  \item \emph{Say it:} \emph{ñāpakam}
  \item \emph{Say it:} \emph{ñāpakam}, once more
  \item \emph{Say it:} \emph{tāttāvukku ippōtu ñāpakam kuṟaintuviṭṭatu.} — Grandfather’s memory has weakened now.
  \item \emph{Recall:} say the Tamil for old age, then the Tamil for an elderly person, then say \emph{ñāpakam} again
\end{itemize}

\begin{tcolorbox}[breakable,colback=teal!4,colframe=teal!35!black,title={Before you move on}]
What does \ta{ஞாபகம்} mean? (\textquotedblleft{}Memory; recollection\textquotedblright{}.) Say it once more.
\end{tcolorbox}
\section[\texorpdfstring{\ta{மங்கு} (maṅku)}{மங்கு (maṅku)}]{\ta{மங்கு} (maṅku) — to fade; to grow dim}

\label{lesson:TA-C504-manku}



\begin{quote}
Before the new one: say the Tamil for an elderly person, then the Tamil for memory.
\end{quote}

\subsection*{You'll want to know: \ta{மங்கு}}
\textbf{\ta{மங்கு}} — \emph{maṅku} — \textquotedblleft{}to fade; to grow dim\textquotedblright{}.

\begin{quote}
\ta{வயதானதால்} \ta{பாட்டியின்} \ta{கண்} \ta{பார்வை} \ta{மங்கியது}.
\end{quote}

\begin{quote}
\emph{vayatāṉatāl pāṭṭiyiṉ kaṇ pārvai maṅkiyatu.}
\end{quote}

\begin{quote}
With age grandmother’s eyesight grew dim.
\end{quote}

Past (it …): \ta{மங்கியது}, \emph{maṅkiyatu}.

Of light, colour and eyesight: the thing fades, so the past is \ta{மங்கியது}.

\begin{cousinweb}
\textbf{\ta{மங்கு}} (\emph{maṅku}) is native Dravidian (DEDR 4750), in the family of \ta{மழுங்கு} \emph{maḻuṅku}, \textquotedblleft{}to grow blunt or dim\textquotedblright{}.
\end{cousinweb}

\begin{grammarlens}[title={What you've built}]
One more word for describing what you lived and felt.
\end{grammarlens}

\subsection*{Your turn}
\begin{itemize}
\raggedright
  \item \emph{Say it:} \emph{maṅku}
  \item \emph{Say it:} \emph{maṅku}, once more
  \item \emph{Say it:} \emph{vayatāṉatāl pāṭṭiyiṉ kaṇ pārvai maṅkiyatu.} — With age grandmother’s eyesight grew dim.
  \item \emph{Recall:} say the Tamil for an elderly person, then the Tamil for memory, then say \emph{maṅku} again
\end{itemize}

\begin{tcolorbox}[breakable,colback=teal!4,colframe=teal!35!black,title={Before you move on}]
What does \ta{மங்கு} mean? (\textquotedblleft{}To fade; to grow dim\textquotedblright{}.) Say it once more.
\end{tcolorbox}
\section[\texorpdfstring{\ta{சுருங்கு} (curuṅku)}{சுருங்கு (curuṅku)}]{\ta{சுருங்கு} (curuṅku) — to shrink; to wrinkle}

\label{lesson:TA-C504-curunku}



\begin{quote}
Before the new one: say the Tamil for memory, then the Tamil for to fade.
\end{quote}

\subsection*{You'll want to know: \ta{சுருங்கு}}
\textbf{\ta{சுருங்கு}} — \emph{curuṅku} — \textquotedblleft{}to shrink; to wrinkle\textquotedblright{}.

\begin{quote}
\ta{வயதாக} \ta{வயதாக} \ta{அவருடைய} \ta{முகத்} \ta{தோல்} \ta{சுருங்கியது}.
\end{quote}

\begin{quote}
\emph{vayatāka vayatāka avaruṭaiya mukat tōl curuṅkiyatu.}
\end{quote}

\begin{quote}
As he grew older, the skin of his face wrinkled.
\end{quote}

Past (it …): \ta{சுருங்கியது}, \emph{curuṅkiyatu}.

Of cloth in the wash too: \ta{சட்டை} \ta{சுருங்கியது}, the shirt shrank.

\begin{cousinweb}
\textbf{\ta{சுருங்கு}} (\emph{curuṅku}) is native Dravidian (DEDR 2687), to shrink, shrivel and wrinkle.
\end{cousinweb}

\begin{grammarlens}[title={What you've built}]
One more word for describing what you lived and felt.
\end{grammarlens}

\subsection*{Your turn}
\begin{itemize}
\raggedright
  \item \emph{Say it:} \emph{curuṅku}
  \item \emph{Say it:} \emph{curuṅku}, once more
  \item \emph{Say it:} \emph{vayatāka vayatāka avaruṭaiya mukat tōl curuṅkiyatu.} — As he grew older, the skin of his face wrinkled.
  \item \emph{Recall:} say the Tamil for memory, then the Tamil for to fade, then say \emph{curuṅku} again
\end{itemize}

\begin{tcolorbox}[breakable,colback=teal!4,colframe=teal!35!black,title={Before you move on}]
What does \ta{சுருங்கு} mean? (\textquotedblleft{}To shrink; to wrinkle\textquotedblright{}.) Say it once more.
\end{tcolorbox}
